\subsection{Theoretical Analysis of SubjectMix: Convex-Hull Invariance and Subject Disentanglement}
\label{sec:ana_subjectmix}

The combined effect of (i) the one-sided SubjectMix augmentation can be made
precise under a mild local-linearity assumption on the encoder and (ii) the multi-positive contrastive loss with frozen
image targets. The
analysis below shows that, at any minimizer of the loss, the
subject-dependent component of the embedding is forced to collapse to a
constant, while the stimulus-dependent component is rigidly anchored to
the image manifold and therefore preserved.

Under the generative model \eqref{eq:genmodel} we write, in a
neighborhood of the training distribution,
\begin{equation}\label{eq:decomp}
    E_\theta(\xv_{s,i})=\Phi(\zv_i)+\Psi(\mathbf{u}_s)+\xiv_{s,i},
\end{equation}
where $\Phi: \Zc \to\R^{d}$ captures the
stimulus-dependent content, $\Psi:\Uc\to\R^{d}$ captures
subject-specific bias, and $\xiv_{s,i}$ is a zero-mean residual absorbing higher-order interactions and measurement noise. This decomposition is the first-order Taylor expansion of $E_\theta\circ g$ around the population subject
$\bar{\uv}$, and is exact for any encoder if $\Phi$ and $\Psi$
are allowed to depend on the marginal distributions of
$\zv$ and $\uv$.
\paragraph{What multi-positive InfoNCE alone enforces}
The bidirectional multi-positive loss
$\Lc_{\mathrm{ctr}}$ (Eqs.~\eqref{eq:loss1}--\eqref{eq:loss2})
admits, at zero temperature and infinite batch, the population minimizer
\begin{equation}
    E_\theta(\xv_{s,i})\;=\;\vv_i,
    \quad\forall\,s\in\Sc,\;i\in\Ic.
    \label{eq:anchor-orig}
\end{equation}
Combined with Eq.~\eqref{eq:decomp} and absorbing residuals,
\eqref{eq:anchor-orig} forces
$\Psi(\uv_s)=\vv_i-\Phi(\zv_i)$, whose right-hand
side does not depend on $s$. Hence $\Psi$ is constant across subjects,
and the constant can be folded into $\Phi$, giving the disentangled
fixed point $\Psi\equiv\mathbf{0}$, $\Phi(\zv_i)=\vv_i$.
\emph{This is a population statement.} In practice, the empirical loss
only sees the values of $\Psi$ at the finite set
$\{\uv_s\}_{s\in\Sc}$, and any encoder that satisfies
Eq.~\eqref{eq:anchor-orig} pointwise on this discrete set is admissible
- including encoders for which $\Psi$ is highly oscillatory between
training subjects and uncontrolled at any unseen $\uv_{s^\star}$.
The cross-subject probe in Appendix~\ref{app:probes} confirms that, without
SubjectMix, the empirical minimizer does in fact preserve a substantial
amount of subject-specific information.

\paragraph{What SubjectMix adds: a convex-hull invariance.}
SubjectMix replaces the $K$ discrete alignment constraints per stimulus
with a continuum of constraints. Recall Eq. \ref{eq:subjmix},
$$\widetilde{\xv}_{s,s',i}
   =\lambda\xv_{s,i}+(1-\lambda)\xv_{s',i},\; \lambda\sim\mathrm{Beta}(\alpha_{\mathrm{mix}},\alpha_{\mathrm{mix}}).$$
The temporal/spatial front-end $T_{\theta_t}\!\circ S_{\theta_s}$ is a
stack of linear operators followed by pointwise nonlinearities; for
inputs sharing the same stimulus and differing only in subject identity,
both EEG segments lie in the same operating regime, so a first-order
expansion gives
\begin{equation}
    E_\theta(\widetilde{\xv}_{s,s',i})
    \;=\;
    \lambda\,E_\theta(\xv_{s,i})
    +(1-\lambda)\,E_\theta(\xv_{s',i})
    + {\deltav}_{\lambda,s,s',i},
    \label{eq:mix-linearity}
\end{equation}
with ${\deltav}_{\lambda,s,s',i}$ a curvature term. Treating
$\widetilde{\xv}_{s,s',i}$ as a positive of $\mathbf{v}_i$
imposes
\begin{equation}
    E_\theta(\widetilde{\xv}_{s,s',i})\;=\;\vv_i
    \quad\forall\,(s,s')\in\widetilde{\Sc},\;\forall\lambda
    \in\mathrm{supp}(\mathrm{Beta}(\alpha_{\mathrm{mix}},\alpha_{\mathrm{mix}})).
    \label{eq:anchor-mix}
\end{equation}
Combining \eqref{eq:anchor-orig}--\eqref{eq:anchor-mix} along the chord
$\gamma_{s,s',i}(\lambda)
 =\lambda\xv_{s,i}+(1-\lambda)\xv_{s',i}$,
\begin{equation*}
    E_\theta\bigl(\gamma_{s,s',i}(\lambda)\bigr)\;=\;\vv_i,
    \qquad\text{for almost every }\lambda\in[0,1].
\end{equation*}
Because $\mathrm{Beta}(\alpha_{\mathrm{mix}},\alpha_{\mathrm{mix}})$
charges every open subset of $(0,1)$, we conclude that, on the chord
$\gamma_{s,s',i}$,
\begin{equation}
    \frac{\mathrm{d}}{\mathrm{d}\lambda}\,
    E_\theta\bigl(\gamma_{s,s',i}(\lambda)\bigr)\;=\;\mathbf{0},
    \label{eq:directional}
\end{equation}
i.e., $E_\theta$ is approximately constant along the chord. By the chain
rule,
\begin{equation*}
    \nabla E_\theta\bigl(\gamma_{s,s',i}(\lambda)\bigr)\,
    (\xv_{s',i}-\xv_{s,i})\;\approx\;\mathbf{0},
\end{equation*}
so the directional derivative of $E_\theta$ along the
\emph{subject-difference vector} vanishes throughout the chord.

\paragraph{From chord-wise invariance to convex-hull invariance.}
Iterating Eq.~\eqref{eq:directional} over all pairs
$(s,s')\in\widetilde{\Sc}$ for a given stimulus $i$, the encoder
is constrained to be approximately constant on every chord connecting
two training-subject inputs. Since the training subject inputs span the
finite set $\Xc_i=\{\xv_{s,i}\}_{s\in\Sc}$,
chord-wise invariance is equivalent (to first order) to invariance on
the convex hull
$\mathrm{conv}(\Xc_i)$. Substituting
Eq.~\eqref{eq:decomp} into Eq.~\eqref{eq:directional} shows that
this invariance is carried entirely by $\Psi$: the stimulus component
$\Phi(\mathbf{z}_i)$ is shared across the chord, so chord-wise variation
can only be due to the subject component. The empirical effect of
SubjectMix is therefore to enforce
\begin{equation}
    \Psi\bigl(\uv\bigr)\approx\mathrm{const},
    \quad\forall\,\uv\in\mathrm{conv}\{\uv_s\}_{s\in\Sc},
    \label{eq:hull-invariance}
\end{equation}
a strictly stronger empirical condition than the discrete invariance
$\Psi(\uv_s)=\mathrm{const}$ implied by
Eq.~\eqref{eq:anchor-orig} alone.
\paragraph{Why this matters for cross-subject generalization.}
Eq.~\eqref{eq:hull-invariance} is precisely the condition relevant to a
held-out subject $s^\star$: provided $\mathbf{u}_{s^\star}$ lies near
$\mathrm{conv}\{\mathbf{u}_s\}_{s\in\mathcal{S}}$; a mild assumption
given that all subjects are healthy adults performing the same visual
task; continuity of $\Psi$ then yields
$\Psi(\mathbf{u}_{s^\star})\approx\mathrm{const}$, and consequently
$E_\theta(\mathbf{x}_{s^\star,i})\approx\mathbf{v}_i$. The standard
multi-positive loss provides this guarantee only at the $K$ seen
subject points and is silent everywhere else; SubjectMix extends it to
an entire convex region.

\paragraph{Image representation is preserved by construction.}
The image embeddings $\{\vv_i\}$ are extracted from a frozen
pretrained vision model. SubjectMix only perturbs the EEG side and
never interpolates the target $\vv_i$; the loss therefore demands
invariance, not interpolation. As a consequence, the stimulus axis
$\Phi$ is rigidly anchored to the fixed image geometry, while the
chord-wise constraint
\eqref{eq:hull-invariance} acts \emph{exclusively} on $\Psi$. This
asymmetry distinguishes SubjectMix from vanilla
Mixup~\cite{zhang2018mixup}, in which the simultaneous interpolation of
inputs and labels encourages smooth interpolation of representations
rather than invariance.

\paragraph{Soft (finite-$\tau$) variance interpretation.}
At finite temperature, exact equality is replaced by a soft alignment
penalty. Standard alignment-uniformity bounds for InfoNCE give, up to
constants,
\begin{equation*}
    \Lc_{\mathrm{ctr}}
    \;\geq\;
    \frac{1}{\tau}\,
    \E_{i,\,p,p'\in\mathcal{P}_i}\!
    \bigl\|E_\theta(p)-E_\theta(p')\bigr\|^{2}
    \;-\;\mathcal{U}_\theta,
\end{equation*}
where $\mathcal{P}_i$ is the set of positives for stimulus $i$.
Including SubjectMix samples expands $\Pc_i$ from
$\{\xv_{s,i}\}_{s\in\Sc}$ to
$\{\xv_{s,i}\}\cup\{\widetilde{\xv}_{s,s',i}\}$, replacing
a sum of $K$ pairwise variance terms by an integral over chords. The
penalty therefore controls a continuous functional of $\Psi$ on
$\mathrm{conv}\{\uv_s\}$ rather than its values at the discrete
training subjects, which is the soft-loss analogue of
Eq.~\eqref{eq:hull-invariance}. Empirically, this is consistent with the
slower rise of the subject-identity probe under SubjectMix
(Appendix~\ref{app:probes}).
